\begin{abstract}
Unified models that jointly perform visual understanding and image generation still depend heavily on curated supervision during post-training, typically requiring human annotations, preference labels, or external reward models. We study whether a unified model can improve both capabilities autonomously from unlabeled images only. We propose a self-evolving framework with three LoRA-based roles on a frozen backbone: a Proposer that generates visual questions, a Solver that answers and evaluates, and a Generator that synthesizes images. Training is driven by internal consistency rather than external supervision. To stabilize learning, we introduce Solver Token Entropy (STE), a continuous difficulty signal that remains informative even when sample-level consistency is degenerate. For generation, we use a multi-scale internal assessment combining QA fidelity and cycle-consistent captioning, coupling understanding quality directly to generation optimization. We validate the same framework across three architecturally distinct unified models (BLIP3o, BAGEL, and VARGPT$_{1.1}$). Relative to each base model, our method improves understanding metrics by +2.0 to +3.6 points across reported benchmarks and improves GenEval by +0.03 on all three backbones. These results support unsupervised, model-agnostic self-evolution for unified understanding and generation.
\end{abstract}
